\documentclass[11pt]{article}

\usepackage[margin=1in]{geometry}
\usepackage{amsmath,amssymb,amsthm}
\usepackage{mathtools}
\usepackage{hyperref}
\usepackage{xcolor}
\usepackage{enumitem}

\hypersetup{colorlinks=true, linkcolor=blue, citecolor=blue, urlcolor=blue}

\newcommand{\R}{\mathbb{R}}
\newcommand{\N}{\mathbb{N}}
\newcommand{\Rstar}{\mathbb{R}^{\star}}
\newcommand{\eps}{\varepsilon}
\let\straightepsilon\epsilon
\renewcommand{\epsilon}{\varepsilon}
\newcommand{\st}{\operatorname{st}}
\newcommand{\Var}{\operatorname{Var}}

\setlist[description]{leftmargin=0pt, style=unboxed, font=\normalfont\bfseries}

\title{Twenty Exercises in Probability on Ordinary Intervals\\
\large Densities, $dx=\eps$, and $\omega$ as the Dirac delta}
\author{Exercises for engineers, in the hyperreal field $\Rstar$}
\date{\today}

\begin{document}
\maketitle


\section*{The two rules you need}

Nothing here is counted, and no hyperfinite sample space is built. The sample
space is an ordinary real interval such as $[0,1]$; all of the hyperreal
content sits in the integral. Fix a positive infinitesimal $\eps$, put
$\omega = 1/\eps$, and take the resolution of the line to be $\eps$: a point
$y$ \emph{is} the half-open dot $[y, y+\eps)$. The integral is then the
hyperfinite left-endpoint Riemann sum with $dx = \eps$,
\[
  \int_{[a,b)} f(x)\,dx
  := \sum_{k=0}^{(b-a)\omega-1} f(a+k\eps)\,\eps ,
\]
evaluated in $\Rstar$ and \emph{not} followed by $\st$. A probability is the
integral of a density, $P(E)=\int_E p$. Because a point is one dot, this
specializes to the identity used on nearly every page:
\[
  \boxed{\;P(\{y\}) = p(y)\,\eps\;}
\]
In particular a uniform law on $[0,1)$ has $p \equiv 1$, so hitting an exact
number has probability $\eps > 0$, not zero --- and $\int_{[0,1)} 1\,dx
= \omega\eps = 1$ exactly, so ordinary interval probabilities are untouched.

An atom of mass $a$ is not a second mechanism bolted onto the density. It is
the density \emph{value} $p(y)=a\omega$, since then $P(\{y\})=a\omega\eps=a$.
The Dirac delta is therefore an ordinary function taking the value $\omega$ on
one dot, and $\st$ is applied only where it is explicitly wanted.

\paragraph{The other model.} These exercises never count anything. The
companion curriculum answers the same questions by counting a hyperfinite
sample space $\{0,\dots,\omega-1\}$ instead, and its solutions are
machine-checked in Lean~4.
\begin{center}\small\url{https://files.pannous.com/hyperreal-exercises.pdf}\end{center}

\paragraph{Theory.} The axioms, the model, and the underlying calculus are
\emph{not} repeated here. They are developed in the companion paper,
\href{https://files.pannous.com/hyperreals-light.pdf}{\emph{Hyperreal Numbers: $\eps\cdot\omega=1$ --- An
Algebraic, Computable Introduction to Infinitesimal Calculus}}
(\url{https://files.pannous.com/hyperreals-light.pdf}), with the machine-checked Lean~4 development at
\url{https://github.com/pannous/hyper-lean}.

\paragraph{Readiness labels.} Each exercise is tagged by what the present
Lean~4 formalization already supports.
\begin{description}
\item[Now] the value and its derivation use only single-dot integrals and
$\Rstar$ arithmetic, which the concrete representation already has.
\item[Partial] the closed form is exact and short, but the derivation sums over
the $(b-a)\omega$ dots of an interval, for which there is no index type yet.
\item[Research] additionally needs transcendental densities, a proved
$\st$-compatibility theorem, or a genuine conditional-law construction.
\end{description}


\section*{Exercises}


\subsection*{Exercise 1 --- The probability of hitting an exact number \hfill \normalfont\small[Now]}

\textbf{Problem.} Let $X$ be uniform on $[0,1)$, so $p(x) = 1$. Compute $P(X = y)$
for a specified $y$, prove that it is positive, and give its standard part.

\textbf{Why this framework.} This is the question the whole construction exists to
answer. Classically the answer is $0$, which then has to be explained away as
``probability zero but not impossible''. Here the answer is an ordinary
algebraic value, and the classical $0$ is recovered as its shadow.


\subsection*{Exercise 2 --- Finitely many targets \hfill \normalfont\small[Now]}

\textbf{Problem.} With $X$ uniform on $[0,1)$, let $A$ be a set of $r$ specified
points, $r$ a standard positive natural number. Find $P(A)$ and the ratio
$P(A)/P(\{y\})$.

\textbf{Why this framework.} Both events are classically null, so their shadows
cannot be compared. The algebraic values differ by exactly the factor one
expects.


\subsection*{Exercise 3 --- Intervals are unharmed, and the total is exactly one \hfill \normalfont\small[Partial]}

\textbf{Problem.} For $X$ uniform on $[0,1)$ and standard $0 \le  c < d \le  1$,
compute $P([c,d))$ and verify $P([0,1)) = 1$ exactly.

\textbf{Why this framework.} A reformulation that changed ordinary interval
probabilities would be useless. The point of the $\eps$-resolution is that
it is invisible at finite scale.


\subsection*{Exercise 4 --- Half-open versus closed, and what a normalization costs \hfill \normalfont\small[Partial]}

\textbf{Problem.} Integrate the constant $1$ over the \emph{closed} interval $[0,1]$. Then
find the density of the uniform law on $[0,1]$ and recompute $P(\{0\})$.

\textbf{Why this framework.} Classically $[0,1)$ and $[0,1]$ have the same measure
and the distinction is pedantry. Here they differ by exactly one dot, and the
difference is the size of a point.


\subsection*{Exercise 5 --- Conditioning on an ordinary interval \hfill \normalfont\small[Now]}

\textbf{Problem.} With $X$ uniform on $[0,1)$ and $y$ in $[a,b)\subset [0,1)$,
compute $P(X = y\mid a \le  X < b)$.

\textbf{Why this framework.} Classically the numerator is $0$ and the quotient is
undefined; here the infinitesimal cancels against a finite denominator and
returns the answer one expects, an infinitesimal rescaled by the width.


\subsection*{Exercise 6 --- A nonuniform density is read off at the point \hfill \normalfont\small[Now/Partial]}

\textbf{Problem.} Let $X$ have the triangular density $p(x) = 2x$ on $[0,1)$.
Verify the normalization, compute $P(X = y)$, and compare two points $y1$ and
$y2$.

\textbf{Why this framework.} The classical statement ``all points have probability
zero'' hides that some points are more likely than others. Here the likelihood
ratio of two points is finite, visible, and equal to the density ratio.


\subsection*{Exercise 7 --- An atom is a density value, not a second mechanism \hfill \normalfont\small[Now/Partial]}

\textbf{Problem.} Let $p$ be the density that equals $\omega/2$ on the dot at $0$
and $1/2$ on $(0,1)$. Show that $p$ is a probability density, compute
$P(\{0\})$ and $P(\{y\})$ for $y \neq  0$, and compare their orders.

\textbf{Why this framework.} Classical theory needs two different objects --- a
density plus a discrete point mass --- and a case split in every formula. Here
one density carries both, because an infinite value is available.


\subsection*{Exercise 8 --- The Dirac delta is the derivative of the step, exactly \hfill \normalfont\small[Now]}

\textbf{Problem.} Let $H$ be the Heaviside step, $H(x) = 1$ for $x \ge  0$ and $0$
otherwise. Compute the algebraic derivative $dH(x) = (H(x+\eps)-H(x))/\eps$
at every $x$, and integrate the result over $[-1,1)$.

\textbf{Why this framework.} Classically $dH$ is not a function and the identity
$integral of dH = 1$ needs distribution theory. Here $dH$ is an ordinary
function taking the value $\omega$ on one dot.


\subsection*{Exercise 9 --- Three orders of point probability \hfill \normalfont\small[Now]}

\textbf{Problem.} At a point $y$, consider densities $p(y) = a\cdot \omega$, $p(y) = b$,
and $p(y) = c\cdot \eps$ with standard positive $a,b,c$. Compute $P(\{y\})$ in
each case and order the three results.

\textbf{Why this framework.} The classical shadow assigns $0$ to the last two and
$a$ to the first, collapsing an entire ordered scale into two values.


\subsection*{Exercise 10 --- The dart, with no grid at all \hfill \normalfont\small[Partial/Now]}

\textbf{Problem.} Let $(X,Y)$ be uniform on the unit square, $p = 1$. Compute the
probability of one specified point.

\textbf{Why this framework.} The counting treatment of the dart needs an
$\omega$ by $\omega$ grid convention. Here the same answer comes from doing two
integrations instead of one.


\subsection*{Exercise 11 --- A segment beats a point by exactly $\omega$ \hfill \normalfont\small[Partial]}

\textbf{Problem.} In the same square, compute $P(Y = X)$, the probability that the
two coordinates agree, and compare it with Exercise 10.

\textbf{Why this framework.} Both events are classically null. The framework must
explain why hitting a line is infinitely easier than hitting a point without
appealing to a chosen grid.


\subsection*{Exercise 12 --- A uniform law on the whole line \hfill \normalfont\small[Partial]}

\textbf{Problem.} Take the ambient line to be $[-\omega, \omega)$. Find the uniform
density on it, compute $P(\{y\})$ and $P([a,b))$, and compare $P(\{y\})$ with the
point probability of the uniform law on $[0,1)$.

\textbf{Why this framework.} Classically there is no uniform distribution on $R$ at
all. Here there is one, and its point probability lands in a strictly rarer
order than that of the unit interval.


\subsection*{Exercise 13 --- Exact moments of the uniform law \hfill \normalfont\small[Research]}

\textbf{Problem.} For $X$ uniform on $[0,1)$, compute $E[X]$, $E[X^2]$, and
$\Var(X)$ from the integral, without taking any limit.

\textbf{Why this framework.} Classically the answers are $1/2$ and $1/12$. The
exact values record, in addition, the bias of the integral convention that
produced them.


\subsection*{Exercise 14 --- The integral convention is visible at order $\eps$ \hfill \normalfont\small[Research/Now]}

\textbf{Problem.} For a standard $f$ on $[0,1)$, compare the left-endpoint,
right-endpoint, and midpoint sums. Apply the result to $f(x) = x$.

\textbf{Why this framework.} Classically all three converge to the same number and
the choice is irrelevant. Here they are different hyperreal values, and their
difference is exactly the information the limit destroys.


\subsection*{Exercise 15 --- A mixed law needs no case split \hfill \normalfont\small[Research]}

\textbf{Problem.} With the density of Exercise 7 --- mass $1/2$ at the point $0$ and
$1/2$ spread uniformly --- compute $E[X]$ and the distribution function
$F(t) = P(X < t)$.

\textbf{Why this framework.} Classical treatment of a mixed law splits every
formula into a discrete sum plus an integral. Here $F = integral of p$ with no
case distinction, and the jump appears automatically.


\subsection*{Exercise 16 --- Point probability is not invariant, and should not be \hfill \normalfont\small[Now/Partial]}

\textbf{Problem.} Let $X$ be uniform on $[0,1)$ and $Y = 2X$. Find the density of
$Y$, compute $P(Y = y)$, and reconcile it with $P(X = y/2) = \eps$.

\textbf{Why this framework.} A reader meeting $P(\{y\}) = \eps$ for the first time
usually asks whether $\eps$ is an absolute ``size of a point''. It is not,
and the change-of-variables rule shows exactly what it is relative to.


\subsection*{Exercise 17 --- An exponential law \hfill \normalfont\small[Research]}

\textbf{Problem.} Let $p(x) = \lambda\cdot e^{-\lambda\cdot x}$ on $[0, \omega)$ with standard
$\lambda > 0$. Compute $P(X = y)$ and $P(X \ge  t)$ for standard $t$, and discuss
the normalization.

\textbf{Why this framework.} The tail of an exponential law on an infinite domain
is where the ambient-line convention and the transcendental functions have to
agree; the framework should not need a separate limiting argument.


\subsection*{Exercise 18 --- Bayes with an atom and a density in one formula \hfill \normalfont\small[Research]}

\textbf{Problem.} A parameter $T$ is $0$ with probability $1/2$ and otherwise
uniform on $(0,1)$. Given $T = t$, an observation $Z$ has a standard density
$f(z | t)$. Compute the posterior $P(T = 0\mid Z = z)$.

\textbf{Why this framework.} This is the classical ``mixed prior'' nuisance, where
the discrete and continuous parts of the prior must be handled by different
machinery and the observation $\{Z = z\}$ is itself a null event.


\subsection*{Exercise 19 --- The Borel--Kolmogorov paradox dissolves \hfill \normalfont\small[Research/Partial]}

\textbf{Problem.} For $(X,Y)$ uniform on the unit square, condition on the diagonal
$D = \{Y = X\}$. Compute $P(X\in [a,b)\mid D)$ and explain why the classical
paradox does not arise.

\textbf{Why this framework.} Classically $P(D) = 0$, conditioning on $D$ requires
choosing a parametrization, and different parametrizations give different
answers --- the Borel--Kolmogorov paradox. Here $P(D) = \eps > 0$, so the
elementary definition applies and the answer is unique.


\subsection*{Exercise 20 --- What must be proved before this is a measure theory \hfill \normalfont\small[Research]}

\textbf{Problem.} Identify precisely what is needed to prove: (a) that the standard
part of the hyperfinite integral is the classical Riemann integral, for every
standard Riemann-integrable $f$; (b) that $P$ is finitely additive over arbitrary disjoint unions of dots;
(c) that $P$ is \emph{not} countably additive in the classical sense, and why that
is not a defect.

\textbf{Why this framework.} An honest curriculum has to state where the
reformulation stops being elementary. The framework's own advertising ---
``probability is just an integral'' --- is only justified once these are settled.

\textbf{Solution sketch.} (a) needs either a transfer principle or an explicit
uniform-continuity estimate: the hyperfinite sum differs from every finite
Riemann sum of mesh $\ge  \eps$ by an infinitesimal, which requires control of
the modulus of continuity, exactly the content the classical limit hides.
(b) holds by construction for unions of dots, since the integral is a sum over
disjoint dots; the work is in defining which sets are unions of dots.
(c) a countable union of dots is \emph{not} a hyperfinite sum, so the classical
countable-additivity axiom is not available and must not be assumed. This is
the same boundary documented in $notes/counting/\sigma-algebra-not-needed.md$:
the framework computes exactly on constructed events and makes no claim about
an arbitrary sigma-algebra.


\section*{Where the solutions are}

Worked solutions are in the companion booklet
\href{https://files.pannous.com/hyperreal-integral-exercise-solutions.pdf}{\emph{Twenty Exercises in Probability on Ordinary Intervals --- Solutions}}:
\begin{center}\small\url{https://files.pannous.com/hyperreal-integral-exercise-solutions.pdf}\end{center}
Try the exercises first: the algebra is short, and the point of each one is
what survives that a real-valued shadow would have deleted.


\end{document}
